\section*{Introduction}
Coordination among the cells drive various functions in a multicellular tissue where cellular decision making is influenced by interactions with the environment composed of extrinsic stimuli and other cells. These interactions also known as cell–cell communication (CCC) play a crucial role in tissue development, biological processes and functions \citep{Armingol2021}. Recently emerged single-cell RNA sequencing (scRNA-seq) technologies have made it possible to analyze tissue architecture at the resolution of individual cells \citep{Svensson2018,scDREAMER} and have shown great potential in investigating CCC \citep{Camp2017}. Computational methods have been developed for inferring cell-cell communication and mediating ligand-receptor (LR) pairs from scRNA-seq datasets by utilising ligand-receptor databases \citep{Efremova2020, Jin2021}. These methods including CellChat \citep{Jin2021}, CellphoneDB \citep{Efremova2020}, iCELLNET \citep{Noel2021}, NATMI \citep{Hou2020}, scTensor \citep{ALMET202112}, and SingleCellSignalR \citep{10.1093/nar/gkaa183} mostly rely on defining functions based on the expression of ligand and receptor pair. Beyond ligand-receptor interactions, methods such as NicheNet \citep{Browaeys2020}, Cytotalk \citep{Dimitrov2022}, SoptSC \citep{10.1093/nar/gkz204} and CCCExplorer \citep{CHOI20151187} focused on inferring the effects of ligand–receptor interactions on possible downstream target gene expression in the receiving cells. While these methods have provided biological insights based on single-cell datasets, they often generate false-positives by failing to account for the spatial context of cell type arrangement in the tissue, as the spatial location information is not captured in scRNA-seq data. Retaining the spatial context is of paramount importance in CCC inference as cellular communications are often spatially constrained where signalling pathways are activated when ligands diffuse from sender cells to neighbouring receiver cells in close proximity \citep{ALMET202112}. \par 

Recent breakthroughs in sequencing \citep{slide-seq,CHEN20221777} and imaging-based \citep{Eng2019} spatial transcriptomic technologies have enabled the measurement of gene expression profiles at varying cellular resolutions from single-cells to spots consisting of multiple or fractions of cells while preserving the spatial context information for the cells or spots. Such datasets have been utilized for spatially mapping cells from scRNA-seq data for inferring ligand-target interactions (e.g., SpaOTsc \citep{Cang2020}). Moreover, CCC inference methods that specifically use spatial data have also been developed. Methods such as Giotto \citep{Dries2021} and NCEM \citep{Fischer2023} use a spatial neighbor graph for inferring cell-cell interactions; CellPhoneDB v3 \citep{Garcia-Alonso2021} uses the spatial information to define microenvironments within which interactions are restricted; stLearn \citep{Pham2020.05.31.125658} and SpatialDM \citep{Li2023} rely on the co-expression of ligand and receptor genes in relation to the spatial location diversity; SVCA \citep{ARNOL2019202}, MISTy \citep{Tanevski2022} and COMMOT \citep{Cang2023} use probabilistic model, machine learning model and optimal transport respectively for identifying spatially restricted ligand-receptor interactions. However, most of these methods focus on ligand-receptor interactions, they do not examine the effect of the cell-cell interaction on the downstream target gene expression. Moreover, most of these methods aim to uncover the interacting cell types globally and thus may be less sensitive in identifying local interactions. While a handful of methods including COMMOT, stLearn and SpatialDM can compute ligand-receptor interactions at the resolution of a spatial location, they fail to quantify the activity of ligand-target gene pairs at the same resolution given a specific spatial context with a fine-grained spatial distribution of cell types. Furthermore, the existing methods employing low-resolution spatial datasets (e.g., 10X Visium) cannot quantify cell type-specific interactions from the mixture of cell types present in a spot. Given that the existing computational methods were developed based on low-resolution spatial datasets, they may provide sub-optimal results for the recent single-cell resolution spatial technologies including NanoString CosMx and 10X Xenium. There is a critical requirement for reliable computational approaches that can elucidate the probable influence of the ligands on target gene expressions in the receiving cell in a particular spatial context with unique cell type composition.
% such as CellChat \citep{Jin2021} and CellphoneDB \citep{Efremova2020}, manually curated databases such as iCELLNET \citep{Noel2021}, iTALK \citep{PPR:PPR66576}, NATMI \citep{Hou2020}, scTensor \citep{ALMET202112}, and SingleCellSignalR \citep{10.1093/nar/gkaa183}, or are independent of such databases (e.g., COMUNET \citep{10.1093/bioinformatics/btaa482} and PyMINEr \citep{TYLER20191951}).

% High-resolution single-cell RNA sequencing (scRNA-seq) data as well as spatial transcriptomics that retain information on spatial localization are paramount to studying cell-cell communication. Methods for inferring cell-cell communication have been developed that rely solely on scRNA-seq and utilise pre-existing ligand-receptor databases such as CellChat \citep{Jin2021} and CellphoneDB \citep{Efremova2020}, manually curated databases such as iCELLNET \citep{Noel2021}, iTALK \citep{PPR:PPR66576}, NATMI \citep{Hou2020}, scTensor \citep{ALMET202112}, and SingleCellSignalR \citep{10.1093/nar/gkaa183}, or are independent of such databases (e.g., COMUNET \citep{10.1093/bioinformatics/btaa482} and PyMINEr \citep{TYLER20191951}). Although these methods have provided robust approaches to infer potential ligand-receptor interactions, they lack the spatial context of cell type arrangement in the tissue, which is necessary when inferring cellular interactions. More importantly, cellular communications are often spatially constrained as signalling pathways are activated when ligands diffuse from sender cells to neighbouring receiver cells in close proximity \citep{ALMET202112}. To address this challenge, additional approaches have been developed which  either spatially map previously collected scRNA-seq data such as SpaOTsc \citep{Cang2020}, or make use of highly informative spatial transcriptomic data such as CellPhoneDBv3 \citep{Garcia-Alonso2021}, Cell2Cell \citep{Armingol2022}, Giotto \citep{Dries2021}, MISTy \citep{Tanevski2022}, stLEARN \citep{Pham2020.05.31.125658}, SVCA \citep{ARNOL2019202}, NICHES \citep{10.1093/bioinformatics/btac775} and COMMOT \citep{Cang2023}. However, some of these methods do not take into account the fine-grained spatial distribution of cell types, or prior information regarding plausible ligand-receptor interactions from scRNA-seq data or examine downstream target gene response caused by ligand–receptor binding. On the other hand, methods such as  Cytotalk \citep{Dimitrov2022}, scMLnet \citep{10.1093/bib/bbaa327}, NicheNet \citep{Browaeys2020}, SoptSC \citep{10.1093/nar/gkz204} and CCCExplorer \citep{CHOI20151187} can infer the interactions of ligands with possible downstream targets by making use of bulk data or scRNA-seq data, however, the informative spatial context is still lacking. Very recently, a handful of methods such as COMMOT and SpaOTsc, attempted to investigate the correlation between ligand-receptor interactions and target gene expression, but these approaches were unable to clearly exhibit probable ligand-target relationships over a spatial context. As a result, there is a critical requirement for reliable computational approaches that can offer a unique perspective on the probable spatial influence that ligands can have on known targets in a particular spatial context with unique cell type composition.

To address these challenges, we introduce Renoir (ligand-taRgEt iNteractions acrOss spatIal topogRaphy), an end-to-end computational framework for spatially charting the ligand-target activities and inferring spatial communication niches that harbor specific ligand-target activities and cell type composition. Renoir requires either single-cell resolution spatial transcriptomic data or matched low-resolution spatial transcriptomic data and scRNA-seq data from the same tissue in order to delineate the spatial influence between possible ligand-target pairs curated from an existing database \citep{Browaeys2020}. Therefore, Renoir allows one to determine which ligands lead to the activation of which target genes in a given spatial context. For each curated ligand-target pair, Renoir quantifies a neighborhood activity score for each spatial location of the spatial transcriptomics data by taking into account the cell type abundances, cell type-specific mRNA abundance values of the ligand and target genes, expression of suitable receptors for the ligand in the cell type harboring target gene, cell type-specific gene entropy measure and a cell type-specific mutual information between the ligand-target pairing as obtained from single-cell resolution spatial data or scRNAseq data. The neighborhood activity score computed by Renoir denotes the strength of the ligand-target activity within the context of a spot’s neighbourhood. In addition to mapping potential ligand-target relationships throughout a spatial topology, the neighborhood scores computed by Renoir further enables a wide variety of downstream analyses. Using the neighbourhood scores, Renoir can infer communication niches, which constitute tissue microenvironments harboring specific ligand-target activities pertaining to major co-localized cell types. The neighborhood scores also allow for the spatial mapping of the total activity of gene sets corresponding to well-established signaling pathways as well as inferred denovo through unsupervised clustering. Renoir can also identify the major ligands active in a communication niche and rank them based on their cumulative activity across the target genes expressed in the niche. \par 

Using a wide range of synthetic datasets generated across a variety of tissue types (brain, cancer and intestine) as well as real datasets, we demonstrated that Renoir outperform the state-of-the-art methods in inferring spatial activity of ligand-target pairs and communication niches. To demonstrate the utility of Renoir in uncovering spatial map of ligand-target activities across tissue types, we applied Renoir to three spatial transcriptomic datasets ranging from development to disease, including adult mouse brain, human triple negative breast cancer (TNBC), and human fetal liver, for which Renoir uncovered distinct spatial domains characterized by specific ligand-target activities. Renoir characterized the interactions of regional astrocyte subtypes with other spatially co-localized cell types in the mouse brain, identified several spatially restricted interactions between cancer, stromal and immune cell populations in TNBC, and uncovered hepatocyte-macrophage interactions in developing fetal liver. To further demonstrate Renoir's ability to handle single-cell resolution spatial datasets, we applied it on a NanoString CosMx dataset from hepatocellular carcinoma for which Renoir identified known and novel onco-fetal interactions.  
% from spatial datasets generated across various technologies and hepatocellular carcinoma